\chapter{Un assistant de preuve au quotidien}\label{ch:lean}

Un assistant de preuve vérifie chaque pas d'une démonstration. On s'en sert ici avec
Lean~4 et sa bibliothèque mathlib, mais les règles valent pour d'autres. Ce chapitre ne
décrit pas le langage : il décrit ce qui, autour du langage, rend un projet de
formalisation tenable sur la durée.

\section{Des versions épinglées}

mathlib change chaque jour : des lemmes sont renommés, d'autres disparaissent. Un projet
fixe donc la version exacte de Lean et de mathlib qu'il utilise, dans les fichiers
\fichier{lean-toolchain} et \fichier{lake-manifest.json}. On met à jour quand on le
décide, d'un coup, et l'on répare ce qui casse dans un commit qui ne fait que cela.

\begin{terminal}
# télécharger mathlib déjà compilée
lake exe cache get
# compiler le projet
lake build
# compiler un seul fichier, plus vite
lake env lean Fichier.lean
\end{terminal}

\section{Compter ce qui manque}

Un énoncé écrit mais non démontré porte \texttt{sorry}. Lean l'accepte avec un
avertissement. Le danger est qu'un \texttt{sorry} oublié soutienne silencieusement un
théorème qu'on croit démontré.

\begin{regle}
On vérifie les axiomes d'un théorème avant de le dire démontré.
\end{regle}

\begin{terminal}
#print axioms monTheoreme
-- attendu : [propext, Classical.choice, Quot.sound]
-- si sorryAx apparaît, une partie de la preuve manque
\end{terminal}

Les trois axiomes attendus sont ceux de la bibliothèque ; tout autre nom est un trou. Un
\texttt{sorry} n'est pas une honte : c'est la marque honnête de ce qui reste à faire. Il
devient un problème quand il se cache dans une \emph{définition} dont dépendent ensuite
des théorèmes, parce qu'alors ces théorèmes parlent d'un objet qui n'existe pas.

\section{Distinguer l'énoncé et sa lecture}

Formaliser un énoncé, c'est choisir une lecture. Une preuve en Lean montre que \emph{cette
lecture} est vraie, pas qu'elle est celle de l'auteur. Dans un projet d'édition, la
formalisation trouve parfois un énoncé faux tel qu'il est écrit ; elle ne dit jamais, à elle
seule, ce que l'auteur voulait dire.

\begin{exemple}
Une note affirme qu'un produit fini de modules plats est plat, et que la réciproque vaut.
La formalisation trouve que la bibliothèque contient déjà la première partie ; la seconde
devient l'énoncé cible. Le fichier le dit en commentaire, en une phrase : ce qui est dans
mathlib, ce qui est démontré ici.
\end{exemple}

\section{Contribuer à la bibliothèque}

Un lemme général dont on a eu besoin manque souvent à mathlib. Le proposer est utile, à
condition de suivre les règles de la communauté :
\begin{itemize}
  \item vérifier que le lemme n'existe pas déjà sous un autre nom ;
  \item suivre les conventions de nommage, de style et de documentation ;
  \item discuter d'abord, sur le forum de la communauté, avant d'ajouter une théorie
        nouvelle ;
  \item déclarer l'usage d'un assistant fondé sur un modèle de langage, comme le
        demandent les règles du projet, et ne pas lui faire écrire les échanges avec les
        relecteurs.
\end{itemize}

\begin{warningblock}
Les règles de contribution des bibliothèques communes évoluent, en particulier sur l'usage
de l'intelligence artificielle. On les relit avant chaque contribution.
\end{warningblock}

\section{Feuilles de route}

Certains projets collectifs publient des \emph{feuilles de route} de formalisation : pour
une théorie, la liste des définitions et des théorèmes à écrire, leurs dépendances et les
énoncés cibles, avec un fichier Lean qui les énonce sans les démontrer. Écrire une telle
feuille à partir de ses propres notes oblige à dire précisément ce qui manque et dans quel
ordre ; la faire relire par d'autres est le meilleur test de cette précision.
